%document
\documentclass[10pt]{beamer}
%theme
\usetheme{metropolis}
% packages
\usepackage{color}
\usepackage{listings}
\usepackage[ngerman]{babel}
\usepackage[utf8]{inputenc}
\usepackage{multicol}

\usepackage{hyperref}
\hypersetup{colorlinks=true, linkcolor=blue, urlcolor=red}

% color definitions
\definecolor{mygreen}{rgb}{0,0.6,0}
\definecolor{mygray}{rgb}{0.5,0.5,0.5}
\definecolor{mymauve}{rgb}{0.58,0,0.82}

\lstset{
    backgroundcolor=\color{white},
    % choose the background color;
    % you must add \usepackage{color} or \usepackage{xcolor}
    basicstyle=\footnotesize\ttfamily,
    % the size of the fonts that are used for the code
    breakatwhitespace=false,
    % sets if automatic breaks should only happen at whitespace
    breaklines=true,                 % sets automatic line breaking
    captionpos=b,                    % sets the caption-position to bottom
    commentstyle=\color{mygreen},    % comment style
    % deletekeywords={...},
    % if you want to delete keywords from the given language
    extendedchars=true,
    % lets you use non-ASCII characters;
    % for 8-bits encodings only, does not work with UTF-8
    frame=single,                    % adds a frame around the code
    keepspaces=true,
    % keeps spaces in text,
    % useful for keeping indentation of code
    % (possibly needs columns=flexible)
    keywordstyle=\color{blue},       % keyword style
    % morekeywords={*,...},
    % if you want to add more keywords to the set
    numbers=left,
    % where to put the line-numbers; possible values are (none, left, right)
    numbersep=5pt,
    % how far the line-numbers are from the code
    numberstyle=\tiny\color{mygray},
    % the style that is used for the line-numbers
    rulecolor=\color{black},
    % if not set, the frame-color may be changed on line-breaks
    % within not-black text (e.g. comments (green here))
    stepnumber=1,
    % the step between two line-numbers.
    % If it's 1, each line will be numbered
    stringstyle=\color{mymauve},     % string literal style
    tabsize=4,                       % sets default tabsize to 4 spaces
    % show the filename of files included with \lstinputlisting;
    % also try caption instead of title
    language = C,
	showspaces = false,
	showtabs = false,
	showstringspaces = false,
	escapechar = ,
}

\def\ContinueLineNumber{\lstset{firstnumber=last}}
\def\StartLineAt#1{\lstset{firstnumber=#1}}
\let\numberLineAt\StartLineAt



\newcommand{\codeline}[1]{
	\alert{\texttt{#1}}
}

% This Document contains the information about this course.

% Authors of the slides
\author{
    Lecturers:
	%name1
	Mirko Jantschke,
	%name2
	Pascal Scholz
	\\
	%creaters of the slides
	Created by: Richard Mörbitz, Manuel Thieme
}


% Fancy Logo
\titlegraphic{\hfill\includegraphics[height=1.25cm]{../templates/fsr_logo_cropped}}


\title{C-Lessons}
\subtitle{Variables}
\date{\today}



% the actual document
\begin{document}

\maketitle

%-------------------------------------------------------------------------------------------------

\begin{frame}{Contents}
	\tableofcontents
\end{frame}

%-------------------------------------------------------------------------------------------------

\begin{frame}{Getting started}
	The slides are at \href{https://fsr.github.io/c-lessons/materials.html}{fsr.github.io/c-lessons/materials.html}\\
	\bigskip
	There will be tasks! You can find them at \href{http://fsr.github.io/c-lessons/}{fsr.github.io/c-lessons}\\
	\bigskip
	If you have questions, use the auditorium group: \href{https://auditorium.inf.tu-dresden.de/de/groups/110804109}{https://auditorium.inf.tu-dresden.de/de/groups/110804109}\\
	\bigskip
	In case of big trouble, write an e-mail to your tutor.\\
	\bigskip
	\bigskip
	*** new only for a limited time ***\\
	Hackerspace every foo from bar to foobar in room biz.
\end{frame}

%-------------------------------------------------------------------------------------------------
\section{Bugs}
%-------------------------------------------------------------------------------------------------

\begin{frame}{It's not a bug...}
	There are different kinds of errors.
	\begin{itemize}
		\item Compiletime errors
		\item Runtime errors (\textit{bugs})
	\end{itemize}\ \\\ \\
	\textit{Compiletime errors} are easily handable since the compiler shows you where to fix them.\\\ \\
	\textit{Bugs} on the other hand are harder to find because you have no idea where to look for them.
\end{frame}

%-------------------------------------------------------------------------------------------------

\begin{frame}{... it's a feature.}
	Bugs can appear due to different reasons
	\begin{itemize}
		\item Variable overflow
		\item Division by zero
		\item Infinite loops / recursions
		\item Range excess
		\item Segmentation fault
		\item Dereferencing \textit{NULL pointers}
		\item ...
	\end{itemize}
\end{frame}

%-------------------------------------------------------------------------------------------------

\begin{frame}{The dungeon}
	We prepared a little ASCII dungeon.\\
	You can find it in the repository (\url{https://github.com/fsr/c-slides}) in folder \textit{materials/1\_before/}\\
	\begin{itemize}
		\item Look at the code and try to understand what should happen.
		\item If you find mistakes, please leave them. We'll fix them later.
		\\\ 
		\item Compile it (with \textit{-std=c99}).
		\\\ 
		\item And now run it.
		\pause
		\item Try to fix all the mistakes using compiler flags.
	\end{itemize}
\end{frame}

%-------------------------------------------------------------------------------------------------
\section{GDB}
%-------------------------------------------------------------------------------------------------

\begin{frame}[fragile]{The \textbf{G}NU \textbf{D}e\textbf{B}ugger}
	There are tools helping with bugs, called debuggers. GDB is one of them.\\
	\bigskip
	To use it
	\begin{itemize}
		\item You have to install the package \textit{gdb}\\
		\item You have to compile your program with the \textit{-g} flag
		\begin{lstlisting}[numbers=none]
$ gcc -g main.c
\end{lstlisting}
		\item After that you can start your program with gdb:
		\begin{lstlisting}[numbers=none]
$ gdb a.out
\end{lstlisting}
	\end{itemize}
\end{frame}

%-------------------------------------------------------------------------------------------------

\begin{frame}{Commands}
	\begin{itemize}
		\item If you started gdb without a file you can load it with \textbf{file} \textit{file\_name}.
		\item Use \textbf{r[un]} to execute the program with gdb.\\
		You should begin with that. It will give you further information about the crash.
		\item You can set an arbitrary amount of breakpoints with \textbf{b[reak]} \textit{line\_number} or \textbf{b[reak]} \textit{function\_name}.\\
		Begin with a breakpoint at the point the program crashes.
		\item Print values with \textbf{p[rint]} \textit{identifier}.
		\item Use \textbf{w[atch]} \textit{identifier} to break and print a variable when it's changed.
	\end{itemize}
\end{frame}

%-------------------------------------------------------------------------------------------------

\begin{frame}{Once you're at a breakpoint}
	\begin{itemize}
		\item Use \textbf{n[ext]} to execute the next program line only.
		\item \textbf{s[tep]} executes the next instruction.
		\item You can jump to the next breakpoint with \textbf{c[ontinue]}.
		\item To see how you have come to this point in the program flow, type \textbf{backtrace} or \textbf{bt}.\\
		This shows you all functions you called to come there.
		\\\ \\
		\item By only hitting the \textit{return key}, you repeat the last entered  command.
	\end{itemize}
	
\end{frame}

%-------------------------------------------------------------------------------------------------

\begin{frame}[fragile]{Conditional breakpoints}
After setting a breakpoint, GDB assigns an ID to it.\\
You can use this ID to extend the functionality of that breakpoint.
	\begin{itemize}
		\item \textbf{con[dition]} \textit{breakpoint\_ID expression} adds a condition to your Breakpoint:
		\begin{lstlisting}[numbers=none,language=bash]
(gdb) br 42
Breakpoint 1 at 0xbada55: file main.c, line 42.
(gdb) condition 1 i@=@@=@3
\end{lstlisting}
		\item For string comparison, set the string before comparing with \textbf{strcmp}:
		\begin{lstlisting}
(gdb) br main.c:42
Breakpoint 13 at 0xdeadbeef: file main.c, line 42.
(gdb) set $string_to_compare = "lolwut"
(gdb) cond 13 strcmp ( $stringtocompare, c ) @=@@=@ 0
\end{lstlisting}
	\end{itemize}
\end{frame}

%-------------------------------------------------------------------------------------------------

\begin{frame}{Now it's up to you}
	\begin{itemize}
		\item Find and fix all bugs in the dungeon.\\\ \\
		\begin{tabular}{|l|l|}
			\hline
			\textbf{file} & load program\\\hline
			\textbf{r[un]} & execute program\\\hline
			\textbf{b[reak]} & set breakpoint\\\hline
			\textbf{p[rint]} & print variable\\\hline
			\textbf{w[atch]} & break and print variable when it changes\\\hline
			\textbf{n[ext]} & execute next line and break\\\hline
			\textbf{s[tep]} & execute next instruction and break\\\hline
			\textbf{c[ontinue]} & execute until next breakpoint\\\hline
			\textbf{backtrace} / \textbf{bt} & How did I end up here?\\\hline
		\end{tabular}
	\end{itemize}
\end{frame}

%-------------------------------------------------------------------------------------------------
\end{document}
